\documentclass[conference]{IEEEtran}
\IEEEoverridecommandlockouts
\usepackage{cite}
\usepackage{amsmath,amssymb,amsfonts}
\usepackage{algorithmic}
\usepackage{graphicx}
\usepackage{textcomp}
\usepackage{xcolor}
\usepackage{booktabs}
\usepackage{multirow}
\usepackage{url}

\def\BibTeX{{\rm B\kern-.05em{\sc i\kern-.025em b}\kern-.08em
    T\kern-.1667em\lower.7ex\hbox{E}\kern-.125emX}}

\begin{document}

\title{Patch-MLP-TS: A Channel-Independent, Attention-Free Architecture\\
for Long-Term Time Series Forecasting}

\author{
\IEEEauthorblockN{Vedant Jadhav}
\IEEEauthorblockA{
Department of Artificial Intelligence and Machine Learning \\
Pimpri Chinchwad University \\
Pune, India \\
Email: [add email]}
}

\maketitle

\begin{abstract}
Transformer-based forecasters equipped with patching, such as PatchTST, have become
strong baselines for long-term time series forecasting, but their attention mechanism
adds computational cost that is not always matched by a proportional accuracy gain
over much simpler linear models. We revisit the patch-based forecasting paradigm
without attention and present Patch-MLP-TS, a channel-independent, all-MLP architecture
that combines patch embedding with cross-patch (token-mixing) and per-patch
(feature-mixing) MLP blocks, in the spirit of MLP-Mixer. During development we identify
and correct a subtle but consequential failure mode in a naive patch-plus-MLP design:
without an explicit patch positional embedding and a mechanism to mix information
across patches, a mean-pooled patch-MLP model discards temporal order entirely and
is consistently outperformed by a simpler whole-window baseline. After correcting
this, our model (v2) is statistically competitive with PatchTST and DLinear on
average across four ETT benchmarks (3-seed mean MSE: 0.348 vs.\ 0.350 for PatchTST
and 0.348 for DLinear), shows a consistent advantage at long forecasting horizons on
two of four datasets, and is competitive with PatchTST on the large-channel
Electricity dataset (321 variates) while training 1.5--1.8$\times$ faster and
performing inference up to 1.3$\times$ faster. We report our results with multi-seed
variance, a set of ablations isolating the contribution of channel-independence and
patching, and an efficiency analysis, and we are explicit about where our method does
not win, to give an accurate picture of when a simple, attention-free patch-MLP design
is and is not preferable to attention-based or purely linear alternatives.
\end{abstract}

\begin{IEEEkeywords}
time series forecasting, multi-layer perceptron, patching, channel independence,
efficient architectures, Transformer alternatives
\end{IEEEkeywords}

\section{Introduction}

Long-term time series forecasting (LTSF) underlies applications from energy load
planning to traffic and financial forecasting. Transformer-based architectures
\cite{zhou2021informer,wu2021autoformer,zhou2022fedformer} were long assumed to be
the strongest approach to LTSF due to their capacity to model long-range
dependencies, but Zeng et al.\ \cite{zeng2023transformers} showed that a family of
embarrassingly simple linear models (LTSF-Linear, including DLinear and NLinear)
outperforms these Transformers on most standard benchmarks, questioning whether
attention is necessary for this task at all.

Subsequent work reconciled these findings by revisiting how Transformers are applied
to time series rather than abandoning them. PatchTST \cite{nie2023patchtst}
segments each univariate channel into patches and applies channel-independent
attention over patch tokens, closing much of the gap to (and in several cases
surpassing) linear baselines. Around the same time, MLP-Mixer-style architectures
such as TSMixer \cite{chen2023tsmixer} demonstrated that alternating time-mixing and
feature-mixing MLP layers, with no attention at all, can match Transformer-level
accuracy at substantially lower cost.

This raises a natural question: can the patching idea that helps PatchTST be
combined with a purely MLP backbone, discarding attention entirely, without losing
the accuracy patching provides? A naive answer -- patch-embed each channel
independently, apply a stack of per-patch MLP blocks, and pool -- turns out to fail
badly. We show that such a design, without a positional signal for patches and
without any mechanism to mix information across patches, is consistently and
substantially outperformed by an unpatched baseline with the same MLP capacity. This
failure mode is worth documenting explicitly, since it is easy to overlook: it is not
that patching is harmful in principle, but that a patch representation with no
notion of order and no cross-patch interaction discards the very information
patching is supposed to preserve.

We correct this with two additions -- a learnable patch positional embedding and a
token-mixing MLP layer that mixes across patches, following the MLP-Mixer design
\cite{tolstikhin2021mlpmixer} -- and refer to the corrected model as Patch-MLP-TS
(v2) throughout this paper. We evaluate it against linear baselines (Linear, NLinear,
DLinear) and PatchTST across four ETT datasets and the large-channel Electricity
dataset, with three random seeds on ETT and two on Electricity, and report ablations
isolating the effect of channel-independence, patching, and patch length, alongside
a measured (not theoretical) efficiency comparison against PatchTST.

Our contributions are:
\begin{itemize}
\item We identify and document a failure mode in naive patch-plus-MLP forecasting
architectures -- loss of patch identity and order at the pooling stage -- and show
that fixing it with positional embedding and cross-patch mixing closes a
previously large accuracy gap.
\item We provide a multi-seed benchmark against linear and attention-based baselines
on four ETT datasets and Electricity, reporting mean and standard deviation rather
than single-run numbers.
\item We provide ablations isolating channel-independence (large, consistent
effect) and patching (smaller, horizon-dependent effect), and a patch-length
sensitivity study.
\item We report measured training and inference time against PatchTST, showing a
consistent speed advantage that does not depend on a parameter-count reduction.
\end{itemize}

We emphasize accuracy claims that are supported by the data reported here rather
than a narrative of universal superiority: our model is competitive on average and
wins clearly in specific, identifiable settings, and we report the settings where
it does not win.

\section{Related Work}

\textbf{Linear and MLP-based forecasters.} LTSF-Linear \cite{zeng2023transformers}
showed that a single linear layer per channel, optionally combined with trend/
seasonal decomposition (DLinear) or last-value normalization (NLinear), matches or
beats considerably more complex Transformer forecasters on standard LTSF
benchmarks. TSMixer \cite{chen2023tsmixer} extends this direction with an all-MLP
architecture that alternates time-mixing and feature-mixing operations, in the
style of MLP-Mixer \cite{tolstikhin2021mlpmixer}, and reports accuracy competitive
with state-of-the-art Transformers.

\textbf{Patch-based Transformers.} PatchTST \cite{nie2023patchtst} segments each
channel into overlapping patches, embeds each patch as a token, and applies a
channel-independent Transformer encoder over the resulting patch sequence. Patching
enlarges the effective receptive field per token and reduces the quadratic cost of
attention relative to point-wise tokenization, and channel-independence was shown
to be important for stable performance. iTransformer \cite{liu2024itransformer}
takes a related but distinct approach, tokenizing whole variates rather than
patches and applying attention across variates instead of across time.

\textbf{Our position.} We combine the patch tokenization of PatchTST with the
attention-free, all-MLP mixing philosophy of TSMixer, and specifically study what
goes wrong when the two are combined naively -- a design point that, to our
knowledge, has not been explicitly documented in prior work on patch-based
forecasters.

\section{Method}

\subsection{Problem Setup}

Given a multivariate lookback window $X \in \mathbb{R}^{L \times C}$ with $L$ time
steps and $C$ channels, the forecasting task is to predict the next $H$ steps
$Y \in \mathbb{R}^{H \times C}$. We adopt channel-independence throughout: each of
the $C$ channels is processed independently by a shared-weight model, and the
per-channel outputs are concatenated to form the final prediction.

\subsection{Patch-MLP-TS Architecture}

For a single channel $x \in \mathbb{R}^{L}$, we split the lookback window into
$N = \lfloor L / P \rfloor$ non-overlapping patches of length $P$, giving
$x_p \in \mathbb{R}^{N \times P}$. Each patch is linearly embedded to dimension $d$:
\begin{equation}
e = x_p W_{\text{patch}} + b_{\text{patch}}, \quad e \in \mathbb{R}^{N \times d}.
\end{equation}
A learnable positional embedding $E_{\text{pos}} \in \mathbb{R}^{N \times d}$ is
added: $e \leftarrow e + E_{\text{pos}}$.

The embedded patches are passed through $K$ mixer blocks. Each block $k$ applies,
in sequence:
\begin{align}
t &= \text{LayerNorm}(e) \\
t &= \text{MLP}_{\text{token}}(t^\top)^\top \quad \text{(mix across the $N$ patches)} \\
e &\leftarrow e + t \\
e &\leftarrow e + \text{MLP}_{\text{feat}}(\text{LayerNorm}(e)) \quad \text{(mix across features)}
\end{align}
where $\text{MLP}_{\text{token}}: \mathbb{R}^{N} \to \mathbb{R}^{N}$ operates on the
patch dimension (transposing $e$ so that patches, not features, form the last axis)
and $\text{MLP}_{\text{feat}}: \mathbb{R}^{d} \to \mathbb{R}^{d}$ is a standard
two-layer feed-forward block with GELU activation, matching the time-mixing /
feature-mixing structure of MLP-Mixer \cite{tolstikhin2021mlpmixer}.

After $K$ blocks, a final LayerNorm is applied and all $N$ patch embeddings are
flattened (not mean-pooled) and projected to the forecast horizon:
\begin{equation}
\hat{y} = \text{Flatten}(e) \, W_{\text{head}} + b_{\text{head}}, \quad
\hat{y} \in \mathbb{R}^{H}.
\end{equation}

\subsection{The v1 Failure Mode}

An earlier version of this architecture (v1) omitted both the positional embedding
$E_{\text{pos}}$ and the token-mixing MLP, applying only feature-mixing blocks to
each patch independently and mean-pooling across patches before the head. This
design has no mechanism to distinguish an early patch from a late one, and the
mean-pool at the output discards whatever order-dependent information survived
the independent per-patch processing. Empirically, v1 was outperformed by a
no-patching ablation (a single MLP over the whole lookback window) in every one of
8 tested (dataset, horizon) settings on ETTh1 and ETTh2, by margins as large as
2$\times$ in MSE. Adding the positional embedding and token-mixing MLP (v2) reverses
this: the patched model outperforms the no-patching ablation in 5 of the same 8
settings (Section \ref{sec:ablation}), and average accuracy across the full ETT
benchmark improves from clearly the worst of five compared models (mean rank
4.75/5) to competitive with the strongest baselines. We report this transition
explicitly because the failure mode -- patch tokenization without any way to
recover patch order or mix across patches -- is a natural design one might arrive
at, and its consequences are severe enough to be worth documenting.

\section{Experimental Setup}

\subsection{Datasets}
We evaluate on the four standard ETT datasets (ETTh1, ETTh2: hourly, 7 channels;
ETTm1, ETTm2: 15-minute, 7 channels) \cite{zhou2021informer} and the Electricity
dataset (hourly, 321 channels, UCI ElectricityLoadDiagrams20112014). Data is split
chronologically 70\%/10\%/20\% into train/validation/test, with per-channel
z-score normalization computed on the training split only.

\subsection{Baselines}
We compare against Linear and NLinear (single linear layer, with and without
last-value normalization) and DLinear (trend/seasonal decomposition with two
linear heads) \cite{zeng2023transformers}, and PatchTST
\cite{nie2023patchtst}, using a from-scratch reimplementation with
patch length 16, stride 8, model dimension 128, 8 attention heads, and 2 encoder
layers. We additionally report a lightweight TSMixer \cite{chen2023tsmixer}
reproduction (hidden dimension 64, 2 mixer blocks) as a single-seed, directional
reference point; this reproduction is not tuned to match published TSMixer numbers
and should not be read as a rigorous comparison.

\subsection{Training Details}
All models use lookback length $L=96$ and are evaluated at horizons
$H \in \{96, 192, 336, 720\}$ (ETT) and $H \in \{96, 720\}$ (Electricity, to limit
compute given the 321-channel width). Patch-MLP-TS uses $P=12$, $d=128$, $K=3$
mixer blocks, dropout 0.1. All models are trained with AdamW, learning rate
$10^{-3}$, weight decay $10^{-4}$, batch size 64 (ETT) or 16--128 (Electricity,
depending on model memory footprint), up to 20 epochs with early stopping
(patience 5) on validation MSE, using mixed-precision training. We report test-set
MSE and MAE.

\subsection{Seeds}
ETT results are averaged over 3 seeds (42, 43, 44); Electricity results are
averaged over 2 seeds (42, 43) due to compute constraints from its larger channel
count. We report mean $\pm$ standard deviation throughout and do not present the
2-seed Electricity results as equivalent in rigor to the 3-seed ETT results.

\section{Main Results}

Table \ref{tab:main} reports test MSE (mean $\pm$ std) for all five models across
four ETT datasets and Electricity, at short (96) and long (720) horizons. Figure
\ref{fig:main} shows the full horizon sweep with error bars.

\begin{table*}[t]
\caption{Test MSE (mean $\pm$ std across seeds). Bold indicates the lowest mean MSE
per row. ETT: 3 seeds. Electricity: 2 seeds (marked with $\dagger$).}
\label{tab:main}
\centering
\small
\begin{tabular}{llccccc}
\toprule
Dataset & $H$ & Linear & NLinear & DLinear & PatchTST & Patch-MLP-TS (ours) \\
\midrule
ETTh1 & 96  & 0.4436$\pm$0.0012 & 0.4506$\pm$0.0010 & 0.4403$\pm$0.0015 & \textbf{0.4152$\pm$0.0038} & 0.4195$\pm$0.0069 \\
ETTh1 & 720 & 0.6617$\pm$0.0059 & 0.7107$\pm$0.0005 & 0.6520$\pm$0.0036 & \textbf{0.6464$\pm$0.0060} & 0.6500$\pm$0.0049 \\
ETTh2 & 96  & 0.1726$\pm$0.0012 & 0.1750$\pm$0.0005 & \textbf{0.1688$\pm$0.0003} & 0.1760$\pm$0.0067 & 0.1742$\pm$0.0032 \\
ETTh2 & 720 & 0.2712$\pm$0.0029 & 0.2992$\pm$0.0006 & 0.2664$\pm$0.0023 & 0.2633$\pm$0.0130 & \textbf{0.2536$\pm$0.0066} \\
ETTm1 & 96  & 0.3571$\pm$0.0007 & 0.3625$\pm$0.0021 & \textbf{0.3558$\pm$0.0006} & 0.4048$\pm$0.0263 & 0.4065$\pm$0.0146 \\
ETTm1 & 720 & 0.5511$\pm$0.0021 & 0.5836$\pm$0.0014 & 0.5495$\pm$0.0016 & 0.5518$\pm$0.0122 & \textbf{0.5346$\pm$0.0067} \\
ETTm2 & 96  & 0.1216$\pm$0.0007 & 0.1214$\pm$0.0003 & 0.1206$\pm$0.0003 & \textbf{0.1146$\pm$0.0002} & 0.1171$\pm$0.0013 \\
ETTm2 & 720 & 0.2320$\pm$0.0039 & 0.2382$\pm$0.0003 & 0.2293$\pm$0.0035 & \textbf{0.2276$\pm$0.0049} & 0.2287$\pm$0.0046 \\
Electricity$^\dagger$ & 96  & 0.1959$\pm$0.0000 & 0.1991$\pm$0.0002 & 0.1958$\pm$0.0001 & \textbf{0.1715$\pm$0.0007} & 0.1766$\pm$0.0041 \\
Electricity$^\dagger$ & 720 & 0.2424$\pm$0.0003 & 0.2538$\pm$0.0000 & 0.2417$\pm$0.0003 & \textbf{0.2315$\pm$0.0001} & 0.2371$\pm$0.0029 \\
\bottomrule
\end{tabular}
\end{table*}

\begin{figure*}[t]
\centering
\includegraphics[width=\textwidth]{figure1_main_benchmark.pdf}
\caption{Test MSE vs.\ forecast horizon for all five models across four ETT
datasets and Electricity. Error bars show standard deviation across seeds (3 for
ETT, 2 for Electricity).}
\label{fig:main}
\end{figure*}

Across the 10 (dataset, horizon) settings in Table \ref{tab:main}, PatchTST attains
the lowest mean MSE in 6, DLinear in 2, and Patch-MLP-TS in 2. Averaged over the
eight 3-seed ETT rows, mean MSE is 0.3479 for DLinear, 0.3480 for Patch-MLP-TS, and
0.3500 for PatchTST -- Patch-MLP-TS and DLinear are statistically indistinguishable
on this average, and both are marginally ahead of PatchTST and plain Linear
(0.3514) on ETT overall.

Patch-MLP-TS shows its clearest, most seed-stable advantage at long horizons on two
of the four ETT datasets: ETTh2 at $H=720$ (0.2536 vs.\ 0.2633 for PatchTST and
0.2664 for DLinear) and ETTm1 at $H=720$ (0.5346 vs.\ 0.5518 for PatchTST and 0.5495
for DLinear). Conversely, ETTm1 at $H=96$ is the weakest setting for both patch-based
models, and notably the least stable: Patch-MLP-TS has std 0.0146 and PatchTST has
std 0.0263 at this setting, an order of magnitude larger than DLinear's 0.0006,
while the simple linear baselines win outright. ETTh2 at $H=96$ and both Electricity
horizons are also won by PatchTST or DLinear. We do not have a diagnosis for the
ETTm1 $H=96$ instability and flag it as a direction for future work rather than
explain it away.

On Electricity specifically -- the only dataset with a substantially larger channel
count (321 vs.\ 7 for ETT) -- Patch-MLP-TS is close to PatchTST at both horizons
(0.1766 vs.\ 0.1715 at $H=96$; 0.2371 vs.\ 0.2315 at $H=720$) while, as shown in
Section \ref{sec:efficiency}, training and running inference substantially faster.
We consider this the most practically relevant result in the paper: near-parity
accuracy with an attention-based model on a wide-channel dataset, at lower
wall-clock cost.

We additionally compare against a lightweight, untuned TSMixer reproduction at
$H \in \{96, 720\}$ across all five datasets; Patch-MLP-TS outperforms it in all 10
settings tested. Given that this TSMixer implementation is not tuned to match
published numbers, we present this as a directional data point rather than a
rigorous comparison against the TSMixer architecture as originally proposed.

\section{Ablation Studies}
\label{sec:ablation}

\subsection{Channel-Independence vs.\ Channel-Mixing}

We compare the default channel-independent Patch-MLP-TS against a variant that
mixes all channels into each patch token instead of processing channels
independently (Table \ref{tab:channelmix}, single seed).

\begin{table}[h]
\caption{Channel-independent vs.\ channel-mixing variants (MSE, single seed).}
\label{tab:channelmix}
\centering
\small
\begin{tabular}{llcc}
\toprule
Dataset & $H$ & Channel-Independent & Channel-Mixing \\
\midrule
ETTh1 & 96  & 0.4272 & 0.5566 \\
ETTh1 & 720 & 0.6548 & 0.9187 \\
ETTh2 & 96  & 0.1768 & 0.2600 \\
ETTh2 & 720 & 0.2520 & 0.4162 \\
\bottomrule
\end{tabular}
\end{table}

Channel-independence wins by a wide, consistent margin in every setting tested --
this is the single cleanest and most confident finding in our ablation study, and
matches the conclusion drawn for PatchTST \cite{nie2023patchtst}: independently
processing channels of a multivariate series with shared weights is more effective
than mixing them at the patch-embedding stage, at least for the datasets studied
here.

\subsection{Patching vs.\ No-Patching}

We compare the full patched model against a variant using the same mixer backbone
but with no patch split (the entire 96-step lookback embedded as a single token),
using the same positional-embedding-free MLP stack (Table \ref{tab:patching},
single seed).

\begin{table}[h]
\caption{Patching vs.\ no-patching (MSE, single seed). Bold: better of the pair.}
\label{tab:patching}
\centering
\small
\begin{tabular}{llcc}
\toprule
Dataset & $H$ & Patched & No-Patching \\
\midrule
ETTh1 & 96  & \textbf{0.4178} & 0.4246 \\
ETTh1 & 192 & \textbf{0.4689} & 0.4748 \\
ETTh1 & 336 & 0.5273 & \textbf{0.5154} \\
ETTh1 & 720 & \textbf{0.6498} & 0.6698 \\
ETTh2 & 96  & \textbf{0.1727} & 0.1730 \\
ETTh2 & 192 & \textbf{0.2089} & 0.2187 \\
ETTh2 & 336 & 0.2211 & \textbf{0.2164} \\
ETTh2 & 720 & 0.2849 & \textbf{0.2606} \\
\bottomrule
\end{tabular}
\end{table}

Patching wins 5 of 8 settings and loses 3, with the losses concentrated at longer
horizons on ETTh2. We read this as evidence that patching's benefit is
horizon-dependent rather than universal in this architecture: it helps at
shorter-to-medium horizons and on ETTh1 broadly, but the advantage narrows and
reverses at $H \geq 336$ on ETTh2. We report this without qualification rather than
selectively presenting only the settings where patching wins, since the
horizon-dependence itself is informative about when this design choice is
appropriate. Figure \ref{fig:ablation} shows both ablations together.

\begin{figure}[h]
\centering
\includegraphics[width=\columnwidth]{figure2_ablation.pdf}
\caption{Ablation comparison: full (patched) model, no-patching variant, and
channel-mixing variant. ETTh1 shown at $H \in \{96, 336, 720\}$ (patching/no-patching
ablation was not run at $H=192$ for this dataset); ETTh2 shown at all four horizons.}
\label{fig:ablation}
\end{figure}

\subsection{Patch Length Sensitivity}

Table \ref{tab:patchlen} and Figure \ref{fig:patchlen} report MSE at $H=96$ for
patch lengths $P \in \{8, 12, 16, 24\}$ (single seed; $P=12$ is the default used
elsewhere in this paper).

\begin{table}[h]
\caption{Patch length sensitivity at $H=96$ (MSE, single seed).}
\label{tab:patchlen}
\centering
\small
\begin{tabular}{lcccc}
\toprule
Dataset & $P=8$ & $P=12$ (default) & $P=16$ & $P=24$ \\
\midrule
ETTh1 & 0.4153 & 0.4272 & \textbf{0.4155} & 0.4205 \\
ETTh2 & 0.1693 & 0.1768 & \textbf{0.1673} & 0.1718 \\
\bottomrule
\end{tabular}
\end{table}

\begin{figure}[h]
\centering
\includegraphics[width=\columnwidth]{figure4_patch_sensitivity.pdf}
\caption{MSE at $H=96$ as a function of patch length $P$ on ETTh1 and ETTh2.}
\label{fig:patchlen}
\end{figure}

All four patch lengths land within approximately 0.01 MSE of each other on both
datasets, indicating the model is not highly sensitive to this hyperparameter, but
notably the default $P=12$ used throughout our main benchmark is not the best
choice at $H=96$ on either dataset -- $P=16$ performs marginally better in both
cases. We did not retune and rerun the full benchmark with $P=16$; we report this
as a limitation and a direction for a follow-up tuning pass rather than
retroactively presenting only the best patch length.

\section{Efficiency Analysis}
\label{sec:efficiency}

Table \ref{tab:efficiency} and Figure \ref{fig:efficiency} report parameter count,
measured wall-clock training time per epoch, and measured inference latency for
Patch-MLP-TS and PatchTST at $H=96$, on a small-channel dataset (ETTh1, 7 channels)
and a large-channel dataset (Electricity, 321 channels).

\begin{table}[h]
\caption{Efficiency comparison, Patch-MLP-TS vs.\ PatchTST at $H=96$.}
\label{tab:efficiency}
\centering
\small
\begin{tabular}{llccc}
\toprule
Dataset & Model & Params & Epoch (s) & Inference (ms/batch) \\
\midrule
ETTh1 & Patch-MLP-TS & 498{,}856 & \textbf{1.68} & 3.45 \\
ETTh1 & PatchTST & 416{,}224 & 2.58 & \textbf{3.11} \\
Electricity & Patch-MLP-TS & 498{,}856 & \textbf{92.15} & \textbf{25.64} \\
Electricity & PatchTST & 416{,}224 & 161.88 & 33.65 \\
\bottomrule
\end{tabular}
\end{table}

\begin{figure}[h]
\centering
\includegraphics[width=\columnwidth]{figure3_efficiency.pdf}
\caption{Measured training time per epoch and inference latency, Patch-MLP-TS vs.\
PatchTST, on ETTh1 (7 channels) and Electricity (321 channels).}
\label{fig:efficiency}
\end{figure}

Parameter count is \emph{not} a consistent advantage for Patch-MLP-TS: at $H=96$ it
has more parameters than PatchTST (499K vs.\ 416K), because our flatten-based
output head scales with $N \times d \times H$ rather than PatchTST's
$N \times d_{\text{model}} \times H$ with a smaller effective $N \times d$ product
at this horizon. At $H=720$ the relationship reverses (1.14M vs.\ 1.38M, roughly
17\% fewer parameters for Patch-MLP-TS), because the two models' patch counts and
head dimensions scale differently with horizon. We therefore do not claim a
blanket parameter-efficiency advantage.

What we observe consistently is a wall-clock speed advantage: Patch-MLP-TS trains
1.54$\times$ faster on ETTh1 and 1.76$\times$ faster on Electricity per epoch, and
performs inference 1.31$\times$ faster on Electricity (PatchTST is marginally
faster at inference on ETTh1, 3.11ms vs.\ 3.45ms). This speed advantage holds
despite Patch-MLP-TS having a comparable or larger parameter count, and reflects
the absence of the $O(N^2)$ attention computation rather than a smaller model. We
consider this the more meaningful and more honestly supportable efficiency claim
than a parameter-count argument.

\section{Limitations}

We report these directly rather than omitting them:

\begin{itemize}
\item The Electricity benchmark uses 2 seeds, not 3, due to compute constraints
from its 321-channel width; we do not present its variance estimates as
equivalent in rigor to the 3-seed ETT results.
\item The channel-mixing, no-patching, and patch-length ablations (Tables
\ref{tab:channelmix}--\ref{tab:patchlen}) are each single-seed; we do not have
variance estimates for these ablation numbers.
\item ETTm1 at $H=96$ is markedly unstable for both patch-based models (standard
deviation an order of magnitude larger than the linear baselines at the same
setting); we have not diagnosed the cause and flag it as future work.
\item The default patch length ($P=12$) used throughout the main benchmark is not
the best-performing value at $H=96$ on either ETTh1 or ETTh2 (Table
\ref{tab:patchlen}); a full retuning pass was out of scope for this study.
\item Our TSMixer baseline is a lightweight, untuned reproduction and should not
be read as a faithful reproduction of published TSMixer results.
\item Patch-MLP-TS wins the main benchmark in 2 of 10 (dataset, horizon) settings
and loses in the remaining 8 to either PatchTST or DLinear; we do not claim
general superiority, only average competitiveness and specific, identifiable
wins at long horizons on two datasets and on the large-channel Electricity
dataset in combination with lower wall-clock cost.
\end{itemize}

\section{Conclusion}

We present Patch-MLP-TS, a channel-independent, attention-free patch-based
architecture for long-term time series forecasting, and document a failure mode we
encountered during its development: a naive patch-embedding-plus-MLP design with no
positional signal and no cross-patch mixing is consistently and substantially
outperformed by a simpler unpatched baseline, because it discards all information
about patch order. Correcting this with a positional embedding and an MLP-Mixer-style
token-mixing layer yields a model that is statistically competitive with DLinear and
PatchTST on average across four ETT benchmarks, wins clearly at long horizons on two
of those four datasets, is competitive with PatchTST on a substantially wider
(321-channel) dataset, and trains and infers faster than PatchTST throughout,
without a consistent parameter-count advantage to explain the speed gain. We report
our results with multi-seed variance where compute allowed, are explicit about the
settings and the ablations where our method does not win, and see the
horizon-dependence of patching's benefit -- rather than a single verdict of
``patching helps'' or ``patching hurts'' -- as itself a useful finding for future
work on patch-based, attention-free forecasting architectures.

\begin{thebibliography}{00}

\bibitem{zhou2021informer}
H. Zhou, S. Zhang, J. Peng, S. Zhang, J. Li, H. Xiong, and W. Zhang,
``Informer: Beyond efficient transformer for long sequence time-series
forecasting,'' in \emph{Proc. AAAI}, 2021.

\bibitem{wu2021autoformer}
H. Wu, J. Xu, J. Wang, and M. Long, ``Autoformer: Decomposition transformers
with auto-correlation for long-term series forecasting,'' in \emph{Proc.
NeurIPS}, 2021.

\bibitem{zhou2022fedformer}
T. Zhou, Z. Ma, Q. Wen, X. Wang, L. Sun, and R. Jin, ``FEDformer: Frequency
enhanced decomposed transformer for long-term series forecasting,'' in
\emph{Proc. ICML}, 2022.

\bibitem{zeng2023transformers}
A. Zeng, M. Chen, L. Zhang, and Q. Xu, ``Are transformers effective for time
series forecasting?'' in \emph{Proc. AAAI}, 2023.

\bibitem{nie2023patchtst}
Y. Nie, N. H. Nguyen, P. Sinthong, and J. Kalagnanam, ``A time series is
worth 64 words: Long-term forecasting with transformers,'' in \emph{Proc.
ICLR}, 2023.

\bibitem{chen2023tsmixer}
S.-A. Chen, C.-L. Li, N. Yoder, S. O. Arik, and T. Pfister, ``TSMixer: An
all-MLP architecture for time series forecasting,'' \emph{Trans. Mach.
Learn. Res.}, 2023.

\bibitem{tolstikhin2021mlpmixer}
I. O. Tolstikhin, N. Houlsby, A. Kolesnikov, L. Beyer, X. Zhai, T. Unterthiner,
J. Yung, A. Steiner, D. Keysers, J. Uszkoreit, M. Lucic, and A. Dosovitskiy,
``MLP-Mixer: An all-MLP architecture for vision,'' in \emph{Proc. NeurIPS},
2021.

\bibitem{liu2024itransformer}
Y. Liu, T. Hu, H. Zhang, H. Wu, S. Wang, L. Ma, and M. Long, ``iTransformer:
Inverted transformers are effective for time series forecasting,'' in
\emph{Proc. ICLR}, 2024.

\end{thebibliography}

\end{document}
